\scanpage{153}
% archive sheet number: 152
% L1–4
Notons que
\[
\nu\ge\nu(n)-2,\ 0
\]
ce qui implique qu'on ne peut avoir $\nu=0$ i.e. $G_0^+\cap Z$ que dans les cas disjoints \(\substack{\nu(n)=2,\ n=2p\ldots\\\nu(n)=1\ \text{avec}\ldots}\)
% L5–7
\[
G_0^+\cap Z=0\quad\Longleftrightarrow\quad
\left\{\begin{aligned}
&\left.\begin{aligned}
 n&=2p^e\quad p\text{ premier impair}\\
 n&=p^e\quad p\text{ premier impair}
 \end{aligned}\right|e\ge1\\
&n=\deleted{\unsure{2^e}}4\quad\deleted{\text{\unsure{ou}}}\ (\text{cas \unsure{octaédral|octaèdral}})
\end{aligned}\right.
\tag*{(65)}
\]
\deleted{Mais on fait attention qu'ici, dans le cas où}

% L8–9
Dans le cas $n=2^{\deleted{\unsure{a}}e}$ avec $e\ge3$, on trouve \deleted{$G^+\cap Z\simeq\mathbb F_2$.}
\[
 {}_2Z\simeq\mathbb F_2^2,\qquad G_0^+\cap Z\simeq\mathbb F_2
 \qquad(\text{pour }n=2^e,e\ge3).
\tag*{(65)}
\]

% L10–11
On va encore \deleted{définir}\inserted{introduire} l'application canonique
\[
\vec F\xrightarrow{\vec\tau_F}\deleted{\unsure{(\mathbb Z/n\mathbb Z)^*}}\unsure{\vec\Pi|\vec\pi}
\tag*{(66)}
\]
% L12–15
\unsure{s'inscrivant|s'insérant} dans le diagramme commutatif \inserted{(en fait un carré cartésien)}
\diagram{\[
\begin{tikzcd}[ampersand replacement=\&]
\vec F \arrow[r,"\vec\tau_F"] \arrow[d,"\mathrm{can}"'] \arrow[dr,phantom,"\text{\uline{cart}}"] \& \vec\Pi \arrow[d,"\mathrm{can}"]\\
F\arrow[r,"\tau_F"'] \& \Pi
\end{tikzcd}\tag*{(67)}
\]}
(cf (24) pour le ``type'' des faces)
% L16–21
en notant que toute face \deleted{\unsure{orientée}} \deleted{\unsure{définit $(f,\omega)$|$\tau_F(f,\omega)$}}, $f\in F$, si $i=\tau_F(f)$ est son type, il y a corresp 1-1 entre l'ens $\unsure{\underline\omega_f|\simeq_f|\vec F_f}$ des deux orientations de $f$, et l'ens $\mu(i)$ des deux \unsure{orientations|orientables} de $P_i$ i.e. des deux éléments de $\vec\Pi$ au dessus de $i$. On trouve de $=$ un diagramme commutatif \inserted{cartésien}
% L22–23
\diagram{\[
\begin{tikzcd}[ampersand replacement=\&]
\vec A \arrow[r,"\vec\tau_A"] \arrow[d,"\mathrm{can}"'] \& \vec\Pi\arrow[d,"\mathrm{can}"]\\
A\arrow[r,dash,"\tau_A"'] \& \Pi\ ,
\end{tikzcd}\tag*{(68)}
\]}
